(1) 界面切向电场与磁场连续。对非磁介质平面波有波阻抗 \(\eta_j=\eta_0/n_j\)。解边界线性方程得 s 振幅反射系数 \(r_s=(n_1\cos\theta_i-n_2\cos\theta_t)/(n_1\cos\theta_i+n_2\cos\theta_t)\)，p 振幅系数按以各自 p 单位矢量为基准的约定为 \(r_p=(n_2\cos\theta_i-n_1\cos\theta_t)/(n_2\cos\theta_i+n_1\cos\theta_t)\)。若改变反射 p 单位矢量定义，\(r_p\) 整体变号但反射率不变。

(2) Brewster 条件是 \(r_p=0\)，即 \(n_2\cos\theta_i=n_1\cos\theta_t\)。与 \(n_1\sin\theta_i=n_2\sin\theta_t\) 联立，得到 \(\theta_i+\theta_t=\pi/2\)，进一步 \(\tan\theta_B=n_2/n_1\)。在此入射角 p 偏振反射电场消失，而 s 偏振一般仍有反射，因此有助于偏振分离。

(3) 对任一无损介质透射传播波，时间平均法向 Poynting 通量与 \(n\cos\theta\,|E|^2\) 成正比，故反射率 \(R=|r|^2\)，透射率 \(\mathcal T=(n_2\cos\theta_t)/(n_1\cos\theta_i)|t|^2\)。将边界求出的振幅 \(t\) 代入可得 \(R+\mathcal T=1\)。尤其介质折射率或角度不同不能仅凭 \(|r|^2+|t|^2=1\) 判断守恒。

(4) 当 \(n_1>n_2\)，Snell 定律要求 \(\sin\theta_t=(n_1/n_2)\sin\theta_i\)。临界角 \(\theta_c=\arcsin(n_2/n_1)\)；超过它折射波法向波数成为纯虚数，倏逝场衰减常数为 \(\kappa=k_0\sqrt{n_1^2\sin^2\theta_i-n_2^2}\)，深度方向场幅 \(\propto e^{-\kappa z}\)。理想单界面无吸收时总反射率为一，但界面外仍有有限倏逝场，不能把电场直接设零。
